\section{Introduction}\label{sec:introduction}

The \(k\)-server problem asks how to move \(k\) servers in a metric
space when requests arrive one at a time. A request is a point that
must be occupied by a server; the cost is the distance moved.
The offline optimum knows all future requests, whereas an online
policy must decide from the revealed prefix. We determine the sharp
worst-case order of the randomized competitive ratio.

More precisely, let \((X,d)\) be a metric space and let
\(s=(s_1,\ldots,s_k)\in X^k\) be the initial positions of \(k\)
labeled servers. Repeated positions are allowed. To serve \(r_t\),
a policy chooses one server and moves it to \(r_t\), leaving all
others fixed. A server already at \(r_t\) can be chosen at zero cost.
For a finite request sequence \(\sigma\), let
\(\cost_{A,s}(\sigma)\) be a policy's total movement and let
\(\OPT_{X,d,k,s}(\sigma)\) be the minimum such cost with full
knowledge of \(\sigma\). Requests are fixed independently of the
policy's random choices, the usual oblivious-adversary model.

\begin{theorem}\label{thm:main}
There is an absolute constant \(C<\infty\) such that, for every
integer \(k\ge2\), every metric space \((X,d)\) with at least \(k+1\)
points, and every \(s\in X^k\), there exist a randomized online policy
\(A\) and a finite \(B\ge0\) satisfying
\begin{equation}\label{eq:main}
 \E\,\cost_{A,s}(\sigma)
 \le C\bigl(\log(k+1)\bigr)^2\OPT_{X,d,k,s}(\sigma)+B
\end{equation}
for every finite request sequence \(\sigma\).
The constant \(C\) is independent of every property of the metric.
The number \(B\) may depend on \((X,d),k,s\), but not on the sequence
or its length. The same policy serves all finite sequences and
horizons. If the initial positions are distinct, one may take \(B=0\).
\end{theorem}

There is no finiteness, boundedness, or discreteness assumption on
\(X\), and no computational-efficiency claim for the policy.
Throughout, logarithms are natural.

\subsection{Context and previous work}

Manasse, McGeoch, and Sleator introduced the \(k\)-server problem as a
common framework for online service problems
\cite{MMS88,MMS90}. They proved the deterministic lower bound \(k\)
on every metric with more than \(k\) points and conjectured that it is
attainable. Koutsoupias and Papadimitriou established the
\((2k-1)\)-competitive work function algorithm \cite{KP95}.
A recent preprint of Coester, Koutsoupias, and Zbysiński announces a
\(k\)-competitive analysis of that algorithm \cite{CKZ26}.
The present paper concerns the different asymptotic behavior made
possible by randomization against oblivious requests.

Paging, the case of a uniform metric, gave the basic reason to expect
such an improvement. Fiat, Karp, Luby, McGeoch, Sleator, and Young
proved a harmonic lower bound and a \(2H_k\)-competitive randomized
marking algorithm \cite{FKLMSY91}; McGeoch and Sleator attained the
exact bound \(H_k=\sum_{i=1}^k i^{-1}\) \cite{MS91}.
For weighted paging, Bansal, Buchbinder, and Naor obtained the optimal
order \(O(\log k)\) through a fractional primal-dual algorithm and
online rounding \cite{BBN12}. These results formed part of the
evidence for the folklore randomized \(k\)-server conjecture, which
predicted \(O(\log k)\) on every metric.

That conjecture was disproved by Bubeck, Coester, and Rabani
\cite[Theorem~11 and Proposition~8]{BCR23}: there exist
\((k+1)\)-point metrics requiring
competitive ratio \(\Omega(\log^2 k)\), even allowing a fixed additive
constant. Their result must be distinguished from a lower bound
on every metric. The latter developed through the decomposition and
metric Ramsey methods of Karloff, Rabani, and Ravid \cite{KRR94},
Blum, Karloff, Rabani, and Saks \cite{BKRS00}, Bartal, Bollobás,
and Mendel \cite{BBM06}, and Bartal, Linial, Mendel, and Naor
\cite{BLMN05}; Bubeck, Coester, and Rabani ultimately obtained the
sharp universal order \(\Omega(\log k)\)
\cite[Corollary~20]{BCR23}.
Theorem~\ref{thm:main} matches the larger, worst-case lower bound.
It establishes the sharp squared-logarithmic order after the negative
resolution of the original logarithmic conjecture.

The main upper-bound approach passes through hierarchically separated
trees. Bartal introduced probabilistic approximation by these
metrics \cite{Bartal96,Bartal98}, and Fakcharoenphol, Rao, and Talwar
obtained expected distortion \(O(\log n)\) for \(n\)-point metrics
\cite{FRT04}. Recursive allocation methods were developed by
Coté, Meyerson, and Poplawski \cite{CMP08} and by Bansal, Buchbinder,
and Naor \cite{BBN10}. Bansal, Buchbinder, Mądry, and Naor then gave
the first polylogarithmic bound for arbitrary finite metrics,
\(O(\log^2 k\,\log^3 n\,\log\log n)\) \cite{BBMN15}.
Bubeck, Cohen, Lee, Lee, and Mądry achieved \(O(\log^2 k)\) on HSTs
using multiscale entropic regularization \cite{BCLLM18}.
The static embedding reduction yields \(O(\log^2 k\log n)\) on
general finite metrics. Its dependence on \(n\) does not give a
bound in \(k\) alone.
Coester and Cosson recently made the HST and finite-metric bounds
polynomial-time, using a total number of random bits independent of the
request length \cite[Theorem~1 and Corollary~2]{CoesterCosson2026}.
Their general finite-metric ratio retains the factor $\log n$.

Changing the embedding as requests arrive is also an established
direction. The dynamic construction of \cite[Section~4]{BCLLM18}
gives \(O(\log^3 k\log(1+\Delta))\) on finite metrics, where
\(\Delta\) is the ratio
of the largest distance to the smallest positive distance.
Lee's evolving, fusible HST framework used heavy-ball nets to merge
clusters near concentrated mass and introduced a fission potential intended
to pay for later splits \cite[Sections~3.1.2--3.2 and~5.3]{Lee18}.
Its general competitive-ratio claim was withdrawn after a gap in
Section~5.3.2 \cite{LeeErratum}. Heavy-ball protection and accounting
for changing partitions are methodological antecedents of our
construction. Here the partition schedule is driven by the posterior
measure of a hidden comparator; it produces a fixed tree of label
strings. We prove the required edit and retained-anchor travel estimates
in Sections~\ref{sec:partitions} and~\ref{sec:travel}, without invoking
a theorem from the withdrawn argument.

Subtree-rank coordinates, internal supermeasures, and
fractional-to-integral transfer already play central roles in
\cite{BBMN15,BCLLM18}. The use of local ball-mass ratios and their
telescoping across scales has a geometric predecessor in measured
descent \cite[Lemma~1.7 and Section~1.3]{KLMN05}.
The distributional estimates are converted into a single policy
by finite minimax and product compactness.


\subsection{Main ingredients}

We first establish a finite distributional bound. Given a known law
of requests, select one offline optimum for each input and use only
its posterior distributions conditional on the revealed prefix.
This does not reveal the realized future.

The construction starts with partitions of the finite metric at several
distance scales. A point's complete tuple of cell labels identifies a
leaf of a fixed tree. The hidden optimum then gives posterior distributions
of the server counts below each tree vertex. We use those distributions to allocate
fractional server mass on the tree, represent that mass by points of
the original metric, and round it to exactly $k$ servers. Changes of
the partitions and of the representatives must be charged along with
movement of mass through the tree. One logarithm comes from the
partition estimate and one from the allocation estimate; rounding
and realization preserve their product.

For exposition we prove the reusable allocation theorem first, then
construct its partition input. The four components are as follows.

\paragraph{Allocation from subtree count distributions.}
On a fixed separated tree, we maintain fractional subtree amounts
with movement \(O(\log(k+1))\) times that of the hidden comparator.
The amounts dominate the serving requirement and are bounded above
by a constant multiple of the posterior expected counts. They may
park at internal vertices, keeping unused mass above children while
their posterior proportions change. The proof uses convex rank functions
and slowly varying proportions. A separate rule for a nearly dominant
child chooses the allocations to the other children first and gives that child
the remaining parent amount, subject to its cap. This fits all children
within their rank bounds without an additional logarithm.
The resulting allocation theorem applies
to general observed hidden trajectories, independently of the later
metric embedding.

\paragraph{A compact variational potential.}
The partition analysis uses a finite-dimensional minimum over
separated auxiliary sites with positive-part capacity constraints.
Concentration forces a nearly full site near the request.
A feasible relocation of that site makes the optimal integrand
have small positive slope there. Ball-mass ratios and inner
moments then telescope over scales, giving a single logarithm in
the measure-motion budget. All geometric radii in this argument
are fixed numerical multiples of the scale.

\paragraph{Partitions adapted to local mass ratios.}
Heavy balls protect highly concentrated requests. Other requests
use chronological rosters in doubling qualification tiers.
Their first-hit loss depends on a local mass ratio rather than on
one worst-case logarithm at every level. Random retirements keep
the rosters bounded; a separate short-lived auxiliary roster pays
the retirement errors. The resulting finite label sets are fixed
independently of the request horizon.

\paragraph{Moving representatives and rounding.}
Fractional mass parked at a tree vertex is represented by an anchor
in the original metric. Retained anchors can move even when no mass
crosses a tree edge. A normalized affine potential and a pointwise
sum of parked-mass coefficients pay this travel without multiplying
the tree movement by another logarithm. A dynamic balanced-rounding
lemma, allowing correlated parent and child states, incurs only an
absolute loss.

These components give two logarithms in total. The final step removes
the finite-metric additive term by independent episodes and reset
blocks, then uses finite-dimensional separation and compactness.
This yields a single policy on all finite inputs of an arbitrary
metric, including unbounded spaces. A matching simulation handles
repeated initial positions.

\subsection{Organization}

Section~\ref{sec:setup} fixes the probabilistic conventions and the
movement accounting. Sections~\ref{sec:allocation},
\ref{sec:proportions}, and~\ref{sec:allocation-finish} prove the
tree allocation theorem. Section~\ref{sec:pilot} develops the
variational potential, and Section~\ref{sec:partitions} constructs
the changing partitions. Sections~\ref{sec:travel}
and~\ref{sec:rounding} control realization and rounding.
Section~\ref{sec:transfer} proves the uniform-policy conclusion.

The geometric and algorithmic estimates are proved below.
The final passage to arbitrary spaces uses the standard compactness
theorem identified in Section~\ref{sec:transfer}; the worst-case lower
bound is the result of Bubeck, Coester, and Rabani cited above.
